\documentclass{article}
\usepackage[T1]{fontenc}
\usepackage{lmodern}
\usepackage{graphicx}
\usepackage{amsmath,amssymb}
\usepackage[a4paper,margin=2cm]{geometry}
\usepackage[absolute,overlay]{textpos}
\setlength{\TPHorizModule}{1mm}
\setlength{\TPVertModule}{1mm}
\usepackage[indonesian]{babel}
\usepackage{hyperref}
\setlength{\emergencystretch}{2em}
% unit: Halaman 55

\begin{document}

% === Halaman 55 (python/native) ===

55

• Setelah panjang kunci diketahui, maka langkah berikutnya 

menentukan huruf-huruf kunci

• Huruf-huruf kunci dapat ditentukan dengan  menggunakan exhaustive 

key search

• Jika panjang kunci = p, maka jumlah kunci yang harus dicoba sampai 

menemukan kunci yang benar adalah adalah maksimal 26p kali.

• Namun lebih sangkil menemukan huruf-huruf kunci dengan 

menggunakan teknik analisis frekuensi.

\end{document}
